%% ma: L11-B14-0022
%% lop: 11
%% chuong: 4
%% bai: 14
%% dang: 11.14.1
%% loai: DS
%% muc: [NB, NB, TH, NB]
%% dap_an: "ĐSSĐ"
%% nguon: "tai-lieu-hoc-tap-mon-toan-lop-11-theo-chuong-trinh-giao-duc-pho-thong-2018.pdf (Lớp Toán Thầy Duy, Phạm Lê Duy, THPT Chu Văn An), Chương 4 Bài 5, Phần 2, Câu 2 (tr.532)"
%% kien_thuc: "Lớp 11 Bài 14 (phép chiếu song song)"
%% trang_thai: chinh-thuc
%% ghi_chu: "Đã duyệt 03/11/2026."
\begin{ex}
Cho hình lập phương $ABCD.A^{\prime}B^{\prime}C^{\prime}D^{\prime}$. Xét phép chiếu song song theo phương chiếu $AA^{\prime}$.
\choiceTF
{\True Hình chiếu của $C^{\prime}$ lên mặt phẳng $(ABCD)$ là $C$.}
{Hình chiếu của $B$ lên mặt phẳng $(A^{\prime}B^{\prime}C^{\prime}D^{\prime})$ là $A^{\prime}$.}
{Gọi $O$ và $O^{\prime}$ lần lượt là tâm của hình vuông $ABCD$ và $A^{\prime}B^{\prime}C^{\prime}D^{\prime}$. Hình chiếu của $O$ lên mặt phẳng $(ABCD)$ là $O^{\prime}$.}
{\True Hình chiếu của $AC$ lên mặt phẳng $(A^{\prime}B^{\prime}C^{\prime}D^{\prime})$ là $A^{\prime}C^{\prime}$.}
\loigiai{\sloppy Phương chiếu $AA^{\prime}\parallel BB^{\prime}\parallel CC^{\prime}\parallel DD^{\prime}$.
a) $CC^{\prime}\parallel AA^{\prime}$ nên $C^{\prime}\mapsto C$. Đúng.
b) $B\mapsto B^{\prime}$, không phải $A^{\prime}$. Sai.
c) $O\in(ABCD)$ nên hình chiếu của $O$ lên $(ABCD)$ là chính $O$, không phải $O^{\prime}$. Sai.
d) $A\mapsto A^{\prime}$, $C\mapsto C^{\prime}$ nên $AC\mapsto A^{\prime}C^{\prime}$. Đúng.}
\end{ex}
